\clearpage
\section*{Introducere}
\addcontentsline{toc}{section}{Introducere}

Declinul vizitelor la bibliotecile și librăriile tradiționale sugerează că
cititorii din zilele noastre doresc să acceseze informații și să citească
conținut fără a vizita personal o bibliotecă. Multe biblioteci și universități
mari au început deja procesul de digitalizare pentru a face materialele
accesibile membrilor și publicului larg. De asemenea, Google a lansat proiectul
Library cu scopul de a ajuta universitățile și bibliotecile să digitalizeze
milioane de cărți. Corporațiile, de diferite dimensiuni, au început, de
asemenea, să analizeze avantajele Bibliotecii Digitale și au început să adopte
biblioteci digitale pentru angajații lor.

Cărțile au o istorie lungă însă au mai multe dezavantaje printre care sunt
făcute din resurse fizice limitate și pot fi obținute mai lent în comparație cu
cărțile digitale sunt doar unele dintre ele.

Modul în care dorim să accesăm
informații și să citim conținut se schimbă rapid, coincizând cu o scădere a
vizitelor tradiționale la bibliotecă.

Pe măsură ce necesitatea accesului la cantități mai mari de informații va
crește, numărul vizitelelor la librăriile tradiționale va scădea, iar
utilizarea cărților electronice va crește. Mai jos sunt enumerate cele mai importante avantaje ale librăriilor online\cite{Benefits}:
\begin{itemize}
\item \textbf{fără limită fizică}: utilizatorul unei biblioteci digitale nu trebuie să meargă fizic la bibliotecă, oameni din întreaga lume ar putea avea acces la aceleași informații, atâta timp cât este disponibilă o conexiune la internet;
\item \textbf{disponibilitate non-stop}: bibliotecile digitale pot fi accesate în orice moment, 24 de ore pe zi și 365 de zile pe an;
\item \textbf{accesări multiple}: aceleași resurse pot fi utilizate în același timp de un număr de utilizatori;
\item \textbf{abordare structurată}: biblioteca digitală oferă acces la conținut mult mai bogat într-o manieră mai structurată(cu excepția librăriei UTM care este imposibil de navigat), adică putem trece cu ușurință de la catalog la cartea respectivă, apoi la un anumit capitol și așa mai departe;
\item \textbf{regăsire informații}: utilizatorul este capabil de a utiliza orice termen de căutare;
\item \textbf{eficiență}: o copie exactă a originalului poate fi făcută de orice număr de ori fără nici o degradare a calității;
\item \textbf{spațiu}: în timp ce bibliotecile tradiționale sunt limitate de spațiul de stocare, bibliotecile digitale au potențialul de a stoca mult mai multe informații, pur și simplu pentru că informațiile digitale necesită foarte puțin spațiu fizic pentru a le conține;
\item \textbf{cost}: costul întreținerii unei biblioteci digitale este mult mai mic decât cel al unei biblioteci tradiționale. O bibliotecă tradițională trebuie să cheltuiască sume mari de bani plătind pentru personal, întreținerea cărților, chirie și cărți suplimentare. Bibliotecile digitale elimină aceste taxe.
\end{itemize}

Proiectul care va fi prezentat mai jos, este o soluției pentru modelarea unei
librării online. Toate beneficiile enumerate ami sus vor fi prezente în acest
sistem. Deși cărțile electronice sunt foarte utile și ușor de accesat ele nu
pot înlocui în totalitate cărțile fizice astfel că și bibliotecile tradiționale
au locul lor.
